% !TeX spellcheck = de_DE

Zur Umsetzung der Überarbeitung wird ein feedbackbasiertes Re-Prompting eingesetzt.
Anstatt einen Text vollständig neu zu erzeugen erhält das LLM Informationen über die bei der Bewertung festgestellten Schwachstellen.
Das Feedback wird in den Prompt integriert und dient als Grundlage für die nächste Iteration des Textes.

Dieser Ansatz verfolgt das Ziel, gezielte Verbesserungen einzelner Qualitätsaspekte zu ermöglichen, ohne bereits gelungene Aspekte unnötig zu verändern.
Gleichzeitig erlaubt das iterative Vorgehen, die Entwicklung der Kennzahlen über mehrere Überarbeitungszyklen zu beobachten und zu analysieren.